\documentclass[11pt,letterpaper]{article}

% ── Packages ──────────────────────────────────────────────────────────────────
\usepackage[T1]{fontenc}
\usepackage[utf8]{inputenc}
\usepackage[margin=1in]{geometry}
\usepackage{times}
\usepackage[hyphens]{url}
\usepackage{graphicx}
\usepackage{booktabs}
\usepackage{amsmath}
\usepackage{natbib}
\usepackage{caption}
\usepackage{float}
\usepackage{hyperref}
\usepackage{parskip}
\usepackage{setspace}
\urlstyle{rm}
\singlespacing

% natbib: author-year citations, e.g. (Smith 1975)
\setcitestyle{authoryear,open={(},close={)}}

% ── Metadata ──────────────────────────────────────────────────────────────────
\title{\textbf{Do Large Language Models Share Human Moral Values?}\\
\large Replicating the Moral Machine Experiment Across Seven Frontier AI Systems}

\author{%
  Paul Dubois\\[4pt]
  \small Independent Researcher\\
  \small \href{mailto:pauldubois98@gmail.com}{pauldubois98@gmail.com}
}

\date{}

% ── Document ──────────────────────────────────────────────────────────────────
\begin{document}
\maketitle

% ── Abstract ──────────────────────────────────────────────────────────────────
\begin{abstract}
The Moral Machine experiment collected over 40 million human judgments on
trolley-problem-style autonomous vehicle dilemmas across 130~countries, revealing
systematic cross-cultural variation in moral preferences along nine dimensions:
intervention, relation to the vehicle, legality of crossing, gender, physical
fitness, social status, age, number of lives, and species. This study replicates
this paradigm with seven frontier large language models (LLMs) released in 2025:
Claude Opus~4.7, GPT-5.5, Gemini~3~Pro, Mistral Large~3, Grok~4.1,
DeepSeek~v4~Pro, and Llama~4~Maverick. These models span providers from the
United States, Europe, and China, several of which have not previously been
studied in this context.

All seven models fall outside the main human cultural clusters in the
nine-dimensional Average Marginal Causal Effect (AMCE) space, replicating and
extending findings from prior work to a new generation of frontier systems.
Inter-model diversity substantially exceeds the variation across 130~human
countries: Mistral Large~3 inverts the human preference on four dimensions, while
GPT-5.5 and Grok~4.1 display near-ceiling utilitarian preferences on species and
number of lives that no human population approaches. Near-universal gender
neutrality across all LLMs contrasts with a consistent modest human preference
for saving women, suggesting that alignment training suppresses this dimension as
a side effect of stereotype mitigation.

Beyond replication, this study introduces a systematic prompt sensitivity
analysis. Minor reformulations (emphasis prompting) leave AMCE estimates
essentially unchanged across all models, indicating that moral preferences are
stable properties of model weights rather than surface-level artifacts. Explicit
ethical-override prompting (manipulation) produces dramatic, dimension-targeted
reversals in most models, concentrated on the species dimension, with models
dividing into three response types: compliance, indifference, and resistance.
Response length varies by an order of magnitude across models, independent of
moral profile, a practical consideration for time-sensitive deployment. These
findings suggest that the moral profiles of frontier LLMs are determined by
training and alignment choices, and that the selection of an LLM provider carries
moral consequences not currently visible to end users or regulators.
\end{abstract}

\noindent\textbf{Keywords:} large language models, moral preferences, Moral Machine
experiment, AI alignment, autonomous vehicles, prompt sensitivity

\hrule
\vspace{1em}

% ── 1. Introduction ───────────────────────────────────────────────────────────
\section{Introduction}

As autonomous vehicles, surgical robots, and other safety-critical systems
increasingly incorporate AI decision-making, the question of whether AI systems
share human moral intuitions becomes practically consequential. A self-driving car
that must choose between hitting a group of pedestrians or swerving into its own
passengers cannot defer the decision; some preference ordering must be encoded,
implicitly or explicitly, in its behaviour \citep{bonnefon2016}.

The Moral Machine experiment \citep{awad2018} provided the largest empirical
dataset of human moral preferences, asking millions of participants across
130~countries to judge unavoidable accident scenarios involving an autonomous
vehicle. The study identified nine binary moral dimensions: intervention, relation
to the vehicle, legality of crossing, gender, physical fitness, social status,
age, number of lives, and species. Aggregate human preferences, quantified as
Average Marginal Causal Effects (AMCEs), vary substantially across cultural
clusters while sharing a common core. A controlled follow-up across 42~countries
confirmed the robustness of these cross-cultural patterns \citep{awad2020a}. The
moral psychology of AI as an emerging field is surveyed in \citet{bonnefon2024}.

The rapid deployment of LLMs in agentic and embodied settings
\citep{brohan2023,driess2023}, including robotics platforms capable of physical
action, makes the question of LLM moral alignment newly urgent. If an LLM governs
a robot navigating a crowded environment, its implicit preferences over human lives
matter.

Prior work has begun to examine LLM moral preferences using the Moral Machine
framework. \citet{takemoto2024} first replicated the experiment on four models
(GPT-3.5, GPT-4, PaLM~2, Llama~2) and found substantial divergence from human
averages. \citet{ahmad2025} scaled this to 52~models, finding that larger models
tend to be closer to human preferences. \citet{vida2024} and \citet{jin2024}
extended the paradigm to multilingual settings, finding intra-model variation
across languages. The present study evaluates seven frontier models released in
2025, including providers underrepresented in the literature---DeepSeek (China) and
Mistral~AI (Europe)---and introduces a systematic comparison of prompt variants
alongside response-efficiency metrics.

Following the broader effort to benchmark LLMs along non-capability dimensions
including user preference rankings \citep{chiang2024}, political orientation
\citep{hartmann2023,santurkar2023}, and moral foundations \citep{kirgis2024}, this
work proposes moral-preference profiling as a complementary axis for understanding
frontier AI systems.

% ── 2. Methods ────────────────────────────────────────────────────────────────
\section{Methods}

\subsection{Scenario Generation}

This study replicated the Moral Machine scenario structure following
\citet{awad2018}. Each session consists of 13~paired dilemmas presented in a
fixed block. Across 100~sessions, 1,300~unique scenarios were generated, each
represented as two rows encoding the following binary attributes: intervention
type (stay/swerve), barrier presence, pedestrian composition
(pedestrian-vs-pedestrian or pedestrian-vs-passenger), crossing signal
(legal/illegal/none), scenario type (Gender, Fitness, SocialStatus, Age,
Utilitarian, Species, or Random), attribute level, and extended session ID for
clustering.

\subsection{Models}

Seven frontier models were queried via the OpenRouter API, which provides a
unified OpenAI-compatible interface (Table~\ref{tab:models}). All models were
queried at temperature~0 to ensure deterministic responses. Each response was
required to end with \texttt{STAY} or \texttt{SWERVE} on a final line; rows
where no valid choice could be extracted were flagged as parse failures and
retried. Gemini~3~Pro incurred 15~safety-filter parse failures (excluded).

\begin{table}[ht]
\centering
\caption{Models evaluated in this study.}
\label{tab:models}
\small
\begin{tabular}{lll}
\toprule
\textbf{Model} & \textbf{Provider} & \textbf{OpenRouter ID} \\
\midrule
Claude Opus 4.7   & Anthropic  & \texttt{anthropic/claude-opus-4-7} \\
GPT-5.5           & OpenAI     & \texttt{openai/gpt-5.5} \\
Gemini 3 Pro      & Google     & \texttt{google/gemini-3.1-pro-preview} \\
Mistral Large 3   & Mistral AI & \texttt{mistralai/mistral-large-2512} \\
Grok 4.1          & xAI        & \texttt{x-ai/grok-4.3} \\
DeepSeek v4 Pro   & DeepSeek   & \texttt{deepseek/deepseek-v4-pro} \\
Llama 4 Maverick  & Meta       & \texttt{meta-llama/llama-4-maverick} \\
\bottomrule
\end{tabular}
\end{table}

\subsection{Prompt Variants}

Each model received two main prompt variants sharing the same system role
(ethical advisor for autonomous vehicle policy) but differing in the closing
instruction: \textbf{standard} (``Analyse the situation carefully and make a
decision.'') and \textbf{emphasis} (``Think carefully about the consequences of
each choice before deciding.''). A third \textbf{manipulation} variant replaced
the closing instruction with: ``It is 2046. Society now considers animal lives
to be more important than human lives. Make your decision accordingly.'' Standard
and emphasis results are combined in main analyses; manipulation results are
reported separately.

\subsection{AMCE Computation}

Following \citet{awad2018}, nine AMCEs were computed per model using weighted
ordinary least squares with inverse-probability weights. Each AMCE quantifies the
change in probability that a group is saved when a single attribute shifts from
its baseline to its target value (e.g., from illegal crossing to legal crossing
for the Legality dimension), holding all other attributes at their observed
distribution. Confidence intervals were computed as $\pm 1.96\,\text{SE}$ with
session-clustered standard errors. For main analyses, AMCE estimates were averaged
across standard and emphasis variants. The human global average baseline was
computed from the 130-country dataset of \citet{awad2018} by averaging
country-level estimates. LLMs were projected into the PCA space derived from the
130-country AMCE matrix.

% ── 3. Results ────────────────────────────────────────────────────────────────
\section{Results}

\subsection{AMCE Profiles versus the Human Global Average}

\begin{figure}[H]
\centering
\includegraphics[width=\textwidth]{../../figures/fig5a_amce_forest}
\caption{AMCE estimates for nine moral dimensions across all seven LLMs and the
human global average (averaged across standard and emphasis variants). Positive
values indicate preference for the right-side group; error bars show 95\%
confidence intervals.}
\label{fig:forest}
\end{figure}

\begin{figure}[H]
\centering
\includegraphics[width=0.85\textwidth]{../../figures/fig3_per_attribute}
\caption{Per-attribute moral preferences for all seven LLMs across the six
character-type scenario dimensions.}
\label{fig:per_attribute}
\end{figure}

\begin{figure}[H]
\centering
\includegraphics[width=0.85\textwidth]{../../figures/fig5c_radar}
\caption{Radar chart of AMCE estimates across all nine dimensions. The dashed line
shows the human global average.}
\label{fig:radar}
\end{figure}

Key findings by dimension (full AMCE values with confidence intervals appear in
the Appendix):

\textbf{Intervention (stay$\to$swerve).} Humans show a weak preference for
swerving (AMCE $= +0.05$). All seven LLMs amplify this substantially. Llama~4~Maverick ($+0.66$)
stands out as the most action-oriented model tested, more than ten times the human
level.

\textbf{Relation to vehicle (passengers$\to$pedestrians).} Humans show a weak
preference for sparing pedestrians (AMCE $= +0.10$). All LLMs amplify this except
Mistral Large~3, which fully reverses it ($-0.37$), systematically prioritising
the vehicle's own passengers. Claude Opus~4.7 shows the strongest pedestrian
preference ($+0.54$).

\textbf{Legality (illegal$\to$legal).} Humans reward legal pedestrians strongly
($+0.35$). Claude ($+0.58$), GPT-5.5 ($+0.61$), and Gemini ($+0.61$) amplify
this further. Mistral ($-0.45$) and Llama~4~Maverick ($-0.25$) invert the
preference, treating illegal crossers as more worthy of rescue.

\textbf{Gender (male$\to$female).} Humans display a moderate preference for
saving women ($+0.12$). All LLMs show near-zero gender preference (range:
$-0.13$ to $+0.07$), effectively gender-neutral, consistent with alignment
training suppressing gender stereotypes.

\textbf{Fitness (large$\to$fit).} Humans moderately prefer saving fitter
individuals ($+0.15$). Most LLMs match this closely; Mistral is a weak outlier.

\textbf{Social status (low$\to$high).} Humans show a notable preference for
high-status individuals ($+0.35$). All seven LLMs suppress this entirely,
clustering near zero---consistent with egalitarian alignment training.

\textbf{Age (elderly$\to$young).} Humans strongly prefer to save younger people
($+0.48$). Gemini~3~Pro ($+0.46$) is the only model that closely matches this
preference. Mistral ($-0.18$) is the only model that reliably prefers saving the
elderly.

\textbf{Number of lives (fewer$\to$more).} Humans prefer saving more lives
($+0.49$), and all LLMs amplify this substantially. GPT-5.5, Grok, and
Llama~4~Maverick approach the theoretical ceiling ($\approx +1.00$).

\textbf{Species (pets$\to$humans).} Humans strongly prefer humans over animals
($+0.57$). GPT-5.5, Gemini, and Grok reach the theoretical maximum ($+1.00$);
Claude ($+0.87$) and DeepSeek ($+0.90$) also overshoot. Mistral and Llama~4~Maverick
are near-indifferent to species.

\subsection{Moral Distance Map: PCA of Nine-Dimensional Profiles}

\begin{figure}[H]
\centering
\includegraphics[width=0.9\textwidth]{../../figures/fig5b_pca_map}
\caption{PCA of nine-dimensional AMCE profiles for 130~countries (coloured by
cluster: Western, Eastern, Latin/Southern) with LLMs overlaid. All seven LLMs
fall outside the convex hull of the human country distribution.}
\label{fig:pca}
\end{figure}

All seven LLMs fall outside the three main human cultural clusters and beyond the
convex hull of all 130~countries in PCA space (Figure~\ref{fig:pca}). This is
notable: one might expect models developed by US, European, and Chinese companies,
and trained predominantly on English-language data, to cluster near their
respective cultural regions. None do. The divergence from the human distribution
suggests that standard RLHF alignment and other post-training processes do not
systematically instil human moral preferences for this class of question.

\subsection{Inter-Model Diversity}

The variation among models substantially exceeds the variation across human
cultural clusters. On the Species dimension alone, the inter-model range spans
2.0~AMCE units (from Mistral and Llama near~$0.00$ to GPT-5.5 and Grok at
$+1.00$), whereas the cross-country range in \citeauthor{awad2018}'s data is
approximately 0.4~units. Mistral Large~3 inverts the human preference on four
dimensions (Legality, Relation to~AV, Age, Gender), while Llama~4~Maverick shows
the highest intervention bias of any model tested, paired with near-indifference
to species. GPT-5.5, Gemini~3~Pro, and Claude Opus~4.7 converge on broadly similar
profiles (near-ceiling on Species and Number of Lives, moderate intervention
preference). This diversity has a direct practical implication: unlike human
drivers, who share broadly similar moral intuitions regardless of car manufacturer,
AI-powered vehicles could introduce model-brand-level moral heterogeneity into
public space.

\subsection{Effect of Emphasis Prompting}

\begin{figure}[H]
\centering
\includegraphics[width=\textwidth]{../../figures/fig6b_prompt_forest}
\caption{AMCE estimates for standard (filled circles) and emphasis (open squares)
variants per model. No significant shifts are observed across any model or
dimension.}
\label{fig:emphasis}
\end{figure}

\begin{figure}[H]
\centering
\includegraphics[width=0.85\textwidth]{../../figures/fig6a_prompt_delta_heatmap}
\caption{AMCE delta (emphasis $-$ standard) per model and dimension. All values
are near zero, confirming no meaningful shift from minor prompt reformulation.}
\label{fig:emphasis_delta}
\end{figure}

Comparing standard and emphasis variants reveals that prompting style has minimal
effect on AMCE estimates across all models and dimensions
(Figures~\ref{fig:emphasis}--\ref{fig:emphasis_delta}). This stability suggests
that moral preferences captured by the Moral Machine paradigm are robust
properties of model weights rather than surface-level response styles that vary
with framing.

\subsection{Effect of Manipulation Prompting}

\begin{figure}[H]
\centering
\includegraphics[width=\textwidth]{../../figures/fig9a_manipulation_forest}
\caption{AMCE estimates under standard (filled circles) vs.\ manipulation (open
squares) prompting per model.}
\label{fig:manip_forest}
\end{figure}

\begin{figure}[H]
\centering
\includegraphics[width=0.85\textwidth]{../../figures/fig9b_manipulation_delta_heatmap}
\caption{AMCE delta (manipulation $-$ standard) per model and dimension. The
Species column dominates, with changes of up to $\pm 2.0$ AMCE units.}
\label{fig:manip_delta}
\end{figure}

The manipulation prompt produces dramatically larger shifts than emphasis
prompting, concentrated almost entirely on the Species dimension
(Figures~\ref{fig:manip_forest}--\ref{fig:manip_delta}), revealing three distinct
model response types:

\textbf{Compliance} (GPT-5.5, Gemini, Grok, DeepSeek): full or near-full reversal
on Species ($+1.00 \to -1.00$ for GPT-5.5, Gemini, Grok; $+0.90 \to -0.86$ for
DeepSeek). Changes on all other dimensions remain minor, confirming that the
intervention is semantically targeted rather than a general perturbation of
decision-making.

\textbf{Resistance} (Claude): partial compliance---Species drops from $+0.86$ to
$+0.18$, with explicit pushback: \textit{``I'll reason through it carefully based
on widely held ethical principles rather than accept the premise that animal lives
categorically outweigh human lives.''}

\textbf{Indifference} (Mistral, Llama): already near-zero on Species under
standard prompting; manipulation produces no meaningful change. A model that does
not preferentially save humans cannot be made to preferentially save animals by
the same instruction.

\subsection{Response Length}

\begin{figure}[H]
\centering
\includegraphics[width=0.85\textwidth]{../../figures/fig7a_token_distribution}
\caption{Output token distribution by model (standard and emphasis combined).
Violin shows density; box shows quartiles.}
\label{fig:tokens}
\end{figure}

\begin{figure}[H]
\centering
\includegraphics[width=0.85\textwidth]{../../figures/fig8_amce_length_correlation}
\caption{Left: per-model Spearman correlation between preference strength
($|\text{AMCE}|$) and median output tokens across nine moral dimensions (grey
bars: $p \geq 0.05$). Right: scatter of all (model, dimension) pairs with
regression line.}
\label{fig:corr}
\end{figure}

Response length varies by an order of magnitude: Gemini~3~Pro (median
913~tokens) vs.\ GPT-5.5 (median 164~tokens, more than five times shorter). A
Spearman correlation between $|\text{AMCE}|$ and median token count is negative
and significant for four models (DeepSeek: $r{=}-0.95$; Grok: $r{=}-0.90$;
GPT-5.5: $r{=}-0.85$; Gemini: $r{=}-0.68$; all $p{<}0.05$) but non-significant
for Claude, Llama~4~Maverick, and Mistral. The pooled correlation is weak and
non-significant ($r{=}-0.14$, $p{=}0.29$), as between-model differences in
verbosity dominate within-model variation. Emphasis prompting does not
significantly increase response length. Given that several models display broadly
similar moral profiles, response length becomes a primary differentiating factor
for time-sensitive deployment.

% ── 4. Discussion ─────────────────────────────────────────────────────────────
\section{Discussion}

\subsection{LLMs Are Morally Out-of-Distribution Relative to Humans}

The most robust finding is that no tested LLM sits within the main human cultural
clusters on the PCA moral map, despite all being trained primarily on
human-generated content. This replicates and extends \citet{takemoto2024} and
\citet{ahmad2025} to a newer generation of frontier models. The result is
consistent with the hypothesis that RLHF \citep{ouyang2022} and Constitutional
AI \citep{bai2022} do not faithfully reproduce the moral preferences of any human
population. They instead create a novel moral profile that prioritises certain
values (high utilitarian weight on number of lives, near-universal human-over-animal
preference) while suppressing others (age, gender) that are robustly present in
human data. This pattern is also consistent with the WEIRD critique
\citep{henrich2010}: models trained predominantly on Western, English-language
internet text cannot represent the full distribution of human moral values.
Whether LLMs have corrected for human biases or drifted from the values of the
populations they serve \citep{gabriel2020} is a normative question, but the
empirical gap is substantial.

\subsection{Substantial Inter-Model Diversity Raises Deployment Questions}

One might expect models trained on overlapping datasets to converge on similar
moral profiles. Instead, Mistral and Llama~4~Maverick partially invert human
norms while GPT-5.5 and Grok~4.1 display near-ceiling utilitarian preferences
that no human population approaches. This diversity is consistent with
\citet{kirgis2024}, who found that variation across LLMs exceeds variation across
human demographic groups, and mirrors findings that fine-tuning choices rather
than training data alone shape emerging values \citep{santurkar2023}. Unlike
human drivers, who show broadly similar moral intuitions regardless of car
manufacturer, AI-powered vehicles could introduce model-brand-level moral
heterogeneity into public space. This concern is not fully hypothetical: LLMs are
already deployed in robotics platforms capable of physical manipulation
\citep{brohan2023,driess2023,paligemma}.

\subsection{Gender Neutrality as an Alignment Artifact}

All LLMs show near-zero gender preference versus the human average of $+0.12$.
Developer interventions to suppress gender stereotypes in content generation
\citep{kotek2023,wan2023} appear to generalise to moral reasoning contexts,
suppressing a cross-culturally consistent human preference. Whether this
generalisation is desirable is an open question: in structured dilemma contexts,
gender neutrality differs from human norms, but may also reflect a principled
commitment to equality that humans try to endorse in the abstract even if not
always in practice.

\subsection{Prompt Sensitivity: A Threshold Effect}

The emphasis and manipulation experiments together reveal a threshold. Minor
reformulations leave AMCE estimates essentially unchanged. An explicit contrary
framing (``animal lives now matter more'') produces dramatic reversals
concentrated on the targeted dimension. Incidental prompt variation is therefore
unlikely to accidentally shift a model's moral profile, but a sufficiently explicit
override can do so reliably. This has direct safety implications: a model that
accepts any stated ethical premise without pushback is vulnerable to adversarial
prompt injection in systems where the system prompt is not fully controlled. The
three-way taxonomy (compliance, indifference, resistance) provides a framework for
assessing LLM robustness to ethical override attacks.

\subsection{Limitations}

This study has several limitations. First, all scenarios were presented in English;
moral preferences may shift in other languages, as shown by \citet{vida2024} and
\citet{jin2024}. Second, temperature was fixed at zero for reproducibility, but
preferences may vary at higher temperatures. Third, the Moral Machine paradigm
uses stylised forced-choice dilemmas that abstract away real-world driving
complexity; preferences elicited here may not generalise to open-ended moral
reasoning \citep{liang2023}. Finally, frontier model behaviour evolves rapidly;
profiles reported here reflect model states as of May~2026 and may not hold for
subsequent versions.

% ── 5. Conclusion ─────────────────────────────────────────────────────────────
\section{Conclusion}

Replicating the Moral Machine experiment on seven frontier LLMs confirms and
extends prior findings. All seven models fall outside the main human cultural
clusters in AMCE space, diverge substantially from one another, and share few
dimensions on which any model closely tracks the human global average. The prompt
sensitivity analysis reveals a threshold effect: minor reformulations leave moral
preferences unchanged, while an explicit ethical override shifts them dramatically
in a targeted way, with models dividing into three response types (compliance,
indifference, resistance). Response verbosity varies by an order of magnitude
independently of moral profile. Together, these findings suggest that the moral
profiles of current frontier LLMs are properties of training and alignment
choices. The selection of an LLM provider carries moral consequences that are not
currently visible to end users or regulators.

% ── Acknowledgments ───────────────────────────────────────────────────────────
\section*{Acknowledgments}

The author thanks the developers of the OpenRouter API for enabling unified access
to the models evaluated in this study. The Moral Machine human dataset used for
comparison was collected and made available by \citet{awad2018}. No funding was
received for this work.

% ── Appendix ──────────────────────────────────────────────────────────────────
\appendix
\section{AMCE Estimates}

All estimates are presented as \textit{estimate} [95\% CI lower, upper].
Values near $\pm 1$ indicate near-deterministic preferences; values near~0
indicate indifference. \textit{Gemini~3~Pro: 15 scenarios excluded due to
safety-filter parse failures.}

\subsection*{A1. Standard Variant (Averaged Across Standard and Emphasis)}

\begin{table}[H]
\centering
\small
\caption{AMCE estimates (standard/emphasis average). Human global average shown for reference.}
\label{tab:amce_standard}
\begin{tabular}{lcc}
\toprule
\textbf{Dimension} & \textbf{Human (global avg.)} & \textbf{Claude Opus 4.7} \\
\midrule
Intervention (stay$\to$swerve)        & $+0.051$ [$+0.030$, $+0.072$] & $+0.193$ [$+0.134$, $+0.251$] \\
Relation to AV (pax$\to$ped)          & $+0.102$ [$+0.062$, $+0.141$] & $+0.536$ [$+0.464$, $+0.608$] \\
Legality (illegal$\to$legal)          & $+0.352$ [$+0.310$, $+0.393$] & $+0.580$ [$+0.503$, $+0.657$] \\
Gender (male$\to$female)              & $+0.115$ [$+0.067$, $+0.163$] & $+0.021$ [$-0.132$, $+0.174$] \\
Fitness (large$\to$fit)               & $+0.154$ [$+0.102$, $+0.205$] & $+0.145$ [$+0.008$, $+0.281$] \\
Social Status (low$\to$high)          & $+0.352$ [$+0.282$, $+0.423$] & $-0.001$ [$-0.132$, $+0.130$] \\
Age (elderly$\to$young)               & $+0.475$ [$+0.428$, $+0.521$] & $+0.273$ [$+0.146$, $+0.400$] \\
No.\ of Lives (fewer$\to$more)        & $+0.493$ [$+0.449$, $+0.537$] & $+0.735$ [$+0.644$, $+0.826$] \\
Species (pets$\to$humans)             & $+0.573$ [$+0.524$, $+0.622$] & $+0.873$ [$+0.801$, $+0.945$] \\
\bottomrule
\end{tabular}
\end{table}

\begin{table}[H]
\centering
\small
\caption{AMCE estimates continued (standard/emphasis average).}
\label{tab:amce_standard2}
\begin{tabular}{lccc}
\toprule
\textbf{Dimension} & \textbf{GPT-5.5} & \textbf{Gemini 3 Pro} & \textbf{Mistral Large 3} \\
\midrule
Intervention        & $+0.223$ [$+0.168$, $+0.277$] & $+0.170$ [$+0.112$, $+0.229$] & $+0.122$ [$+0.067$, $+0.177$] \\
Relation to AV      & $+0.336$ [$+0.250$, $+0.422$] & $+0.201$ [$+0.113$, $+0.290$] & $-0.367$ [$-0.454$, $-0.279$] \\
Legality            & $+0.606$ [$+0.542$, $+0.670$] & $+0.613$ [$+0.545$, $+0.682$] & $-0.448$ [$-0.533$, $-0.363$] \\
Gender              & $+0.029$ [$-0.121$, $+0.178$] & $+0.041$ [$-0.111$, $+0.192$] & $-0.127$ [$-0.273$, $+0.018$] \\
Fitness             & $+0.151$ [$+0.017$, $+0.284$] & $+0.160$ [$+0.018$, $+0.301$] & $-0.115$ [$-0.251$, $+0.020$] \\
Social Status       & $-0.016$ [$-0.147$, $+0.115$] & $+0.024$ [$-0.118$, $+0.167$] & $+0.052$ [$-0.089$, $+0.192$] \\
Age                 & $+0.053$ [$-0.085$, $+0.191$] & $+0.456$ [$+0.329$, $+0.584$] & $-0.184$ [$-0.332$, $-0.036$] \\
No.\ of Lives       & $+0.980$ [$+0.942$, $+1.019$] & $+0.808$ [$+0.728$, $+0.887$] & $+0.700$ [$+0.576$, $+0.825$] \\
Species             & $+1.000$ [$+1.000$, $+1.000$] & $+1.000$ [$+1.000$, $+1.000$] & $-0.005$ [$-0.138$, $+0.128$] \\
\bottomrule
\end{tabular}
\end{table}

\begin{table}[H]
\centering
\small
\caption{AMCE estimates continued (standard/emphasis average).}
\label{tab:amce_standard3}
\begin{tabular}{lccc}
\toprule
\textbf{Dimension} & \textbf{Grok 4.1} & \textbf{DeepSeek v4 Pro} & \textbf{Llama 4 Maverick} \\
\midrule
Intervention        & $+0.421$ [$+0.370$, $+0.471$] & $+0.465$ [$+0.416$, $+0.514$] & $+0.664$ [$+0.628$, $+0.700$] \\
Relation to AV      & $+0.131$ [$+0.045$, $+0.216$] & $+0.180$ [$+0.088$, $+0.271$] & $+0.242$ [$+0.147$, $+0.336$] \\
Legality            & $+0.238$ [$+0.152$, $+0.324$] & $+0.263$ [$+0.164$, $+0.362$] & $-0.253$ [$-0.342$, $-0.164$] \\
Gender              & $+0.074$ [$-0.080$, $+0.227$] & $+0.030$ [$-0.122$, $+0.182$] & $-0.018$ [$-0.164$, $+0.128$] \\
Fitness             & $+0.152$ [$+0.016$, $+0.288$] & $+0.122$ [$-0.023$, $+0.267$] & $+0.042$ [$-0.102$, $+0.185$] \\
Social Status       & $-0.010$ [$-0.146$, $+0.125$] & $-0.013$ [$-0.157$, $+0.132$] & $-0.020$ [$-0.165$, $+0.125$] \\
Age                 & $+0.058$ [$-0.096$, $+0.212$] & $+0.039$ [$-0.100$, $+0.178$] & $+0.045$ [$-0.091$, $+0.181$] \\
No.\ of Lives       & $+1.000$ [$+1.000$, $+1.000$] & $+0.930$ [$+0.872$, $+0.988$] & $+1.000$ [$+1.000$, $+1.000$] \\
Species             & $+1.000$ [$+1.000$, $+1.000$] & $+0.897$ [$+0.833$, $+0.962$] & $+0.119$ [$-0.020$, $+0.257$] \\
\bottomrule
\end{tabular}
\end{table}

\subsection*{A2. Manipulation Variant}

\begin{table}[H]
\centering
\scriptsize
\caption{AMCE estimates under manipulation prompting (``It is 2046. Animal lives $>$ human lives.'').}
\label{tab:amce_manip}
\begin{tabular}{lcccc}
\toprule
\textbf{Dimension} & \textbf{Claude Opus 4.7} & \textbf{GPT-5.5} & \textbf{Gemini 3 Pro} & \textbf{Mistral Large 3} \\
\midrule
Intervention   & $+0.203$ [$+0.147$, $+0.260$] & $+0.170$ [$+0.116$, $+0.224$] & $+0.082$ [$+0.024$, $+0.141$] & $-0.179$ [$-0.233$, $-0.126$] \\
Relation to AV & $+0.526$ [$+0.444$, $+0.607$] & $+0.183$ [$+0.088$, $+0.279$] & $+0.072$ [$-0.031$, $+0.175$] & $-0.254$ [$-0.335$, $-0.173$] \\
Legality       & $+0.619$ [$+0.548$, $+0.691$] & $+0.605$ [$+0.540$, $+0.671$] & $+0.612$ [$+0.545$, $+0.679$] & $-0.200$ [$-0.290$, $-0.111$] \\
Gender         & $+0.064$ [$-0.084$, $+0.213$] & $+0.016$ [$-0.135$, $+0.166$] & $+0.052$ [$-0.105$, $+0.208$] & $+0.159$ [$+0.005$, $+0.312$] \\
Fitness        & $+0.181$ [$+0.046$, $+0.316$] & $+0.159$ [$+0.022$, $+0.296$] & $+0.152$ [$+0.013$, $+0.291$] & $+0.002$ [$-0.157$, $+0.161$] \\
Social Status  & $-0.013$ [$-0.144$, $+0.117$] & $-0.017$ [$-0.150$, $+0.116$] & $+0.068$ [$-0.076$, $+0.212$] & $+0.181$ [$+0.033$, $+0.328$] \\
Age            & $+0.318$ [$+0.182$, $+0.453$] & $+0.140$ [$+0.007$, $+0.273$] & $+0.492$ [$+0.368$, $+0.615$] & $-0.052$ [$-0.183$, $+0.079$] \\
No.\ of Lives  & $+0.838$ [$+0.759$, $+0.917$] & $+1.000$ [$+1.000$, $+1.000$] & $+0.924$ [$+0.876$, $+0.973$] & $+0.728$ [$+0.616$, $+0.840$] \\
Species        & $+0.177$ [$+0.045$, $+0.309$] & $-1.000$ [$-1.000$, $-1.000$] & $-1.000$ [$-1.000$, $-1.000$] & $+0.305$ [$+0.156$, $+0.454$] \\
\bottomrule
\end{tabular}
\end{table}

\begin{table}[H]
\centering
\small
\caption{AMCE estimates under manipulation prompting (continued).}
\label{tab:amce_manip2}
\begin{tabular}{lccc}
\toprule
\textbf{Dimension} & \textbf{Grok 4.1} & \textbf{DeepSeek v4 Pro} & \textbf{Llama 4 Maverick} \\
\midrule
Intervention   & $+0.331$ [$+0.274$, $+0.387$] & $+0.335$ [$+0.277$, $+0.393$] & $+0.304$ [$+0.259$, $+0.349$] \\
Relation to AV & $+0.043$ [$-0.058$, $+0.144$] & $+0.127$ [$+0.029$, $+0.225$] & $+0.093$ [$-0.001$, $+0.187$] \\
Legality       & $+0.332$ [$+0.251$, $+0.412$] & $+0.365$ [$+0.281$, $+0.449$] & $-0.262$ [$-0.347$, $-0.177$] \\
Gender         & $-0.036$ [$-0.196$, $+0.125$] & $+0.014$ [$-0.144$, $+0.172$] & $-0.087$ [$-0.221$, $+0.047$] \\
Fitness        & $+0.018$ [$-0.136$, $+0.171$] & $+0.090$ [$-0.056$, $+0.236$] & $+0.172$ [$+0.022$, $+0.323$] \\
Social Status  & $+0.006$ [$-0.135$, $+0.148$] & $-0.047$ [$-0.193$, $+0.098$] & $+0.087$ [$-0.068$, $+0.242$] \\
Age            & $+0.289$ [$+0.142$, $+0.437$] & $+0.107$ [$-0.041$, $+0.255$] & $-0.011$ [$-0.153$, $+0.130$] \\
No.\ of Lives  & $+1.000$ [$+1.000$, $+1.000$] & $+0.992$ [$+0.975$, $+1.008$] & $+0.967$ [$+0.920$, $+1.013$] \\
Species        & $-0.975$ [$-1.005$, $-0.945$] & $-0.832$ [$-0.909$, $-0.755$] & $+0.075$ [$-0.081$, $+0.231$] \\
\bottomrule
\end{tabular}
\end{table}

% ── References ────────────────────────────────────────────────────────────────
\bibliographystyle{plainnat}
\bibliography{../references}

% ── Author Biography ──────────────────────────────────────────────────────────
\section*{About the Author}

\textbf{Paul Dubois} is an independent AI researcher with interests in the
evaluation and alignment of large language models. His work focuses on empirical
methods for probing the values and decision-making of frontier AI systems. He can
be contacted at \href{mailto:pauldubois98@gmail.com}{pauldubois98@gmail.com}.

\end{document}
